\section{Corrected source inputs and actual local outputs}
\label{sec:peps_corrected_input_frames}

We construct the local input isometries and output expansion used in
\cite[proof of Theorem~5.2]{OpenAI2026PolynomialPEPS}.
% Source: 04-compression.tex, lines 409--434, at
% adc7f1241b42e322a6451854ab7e4b4c146bf78a.

Source occurrences retain their chronological positions, even when their
endpoint parties coincide with those of another occurrence.
A register list $L$ has memory $\mathcal H_L$; concatenation is identified
with tensor product by the canonical isometry
$\alpha_{L,M}\colon\mathcal H_{L+M}\to\mathcal H_L\otimes\mathcal H_M$.
For a division of the parties into two classes, write $\Pi_L$ for the
isometry gathering the two classes in their inherited orders.
Relabelling parties changes neither the order nor the Hilbert spaces.

\begin{theorem}[Identifying equal register lists]
    \uses{thm:peps_basis_memory_equalities}
    \label{thm:peps_register_identifications}
    \lean{TNLean.PEPS.PairEffect.Word.eval_castLayouts_apply_heq}
    \lean{TNLean.PEPS.PairEffect.Layout.memCongr_append_tmul}
    \lean{TNLean.PEPS.PairEffect.Word.eval_source_apply_heq}
    \leanok
    Identifying two equal register lists leaves their vectors and the evaluation
    of a composition unchanged. Such identifications commute with source
    preparation and with the concatenation isometries $\alpha_{L,M}$.
\end{theorem}
\begin{proof}
    \leanok
    Identify the equal lists. Each assertion is then an equality of the same
    vector or linear map.
\end{proof}

\begin{definition}[Ordered source positions]
    \label{def:peps_ordered_source_positions}
    \lean{TNLean.PEPS.PairEffect.SourceInventory.sourcePositionsLayout}
    \lean{TNLean.PEPS.PairEffect.SourceInventory.prepareSourcePositions}
    \lean{TNLean.PEPS.PairEffect.SourceInventory.slotReference}
    \lean{TNLean.PEPS.PairEffect.SourceInventory.sourcePositionsLayout_eq_layout}
    \lean{TNLean.PEPS.PairEffect.SourceInventory.sourcePositionsLayout_append}
    \lean{TNLean.PEPS.PairEffect.SourceInventory.selectedSources_eq_map}
    \lean{TNLean.PEPS.PairEffect.SourceInventory.sourcePositionsLayout_filter}
    \leanok
    For a family of pair sources $r_i$ and a finite ordered list $J$ of positions,
    let $L_J$ be the concatenation of their two-register lists, and let
    $P_J(\eta)$ prepare the vectors $\eta_i$ in that order.
    Then $L_{J+K}=L_J+L_K$. Selecting positions by a predicate retains their
    inherited order and their literal endpoint parties and spaces. The same
    layout is obtained by assigning zero vectors to these fixed spaces.
\end{definition}

\begin{theorem}[Preparation in an ordered list]
    \label{thm:peps_ordered_source_preparation}
    \lean{TNLean.PEPS.PairEffect.SourceInventory.eval_prepareSourcePositions_heq}
    \lean{TNLean.PEPS.PairEffect.SourceInventory.eval_prepareSourcePositions_append}
    \lean{TNLean.PEPS.PairEffect.SourceInventory.eval_prepareSourcePositions_filter}
    \leanok
    Preparation indexed by positions equals preparation of the corresponding
    ordered source records. Under the concatenation identifications,
    $P_{J+K}(\eta)=P_J(\eta)P_K(\eta)$. Restricting $J$ to a selected set of
    positions agrees with preparing precisely those selected sources.
\end{theorem}
\begin{proof}
    \leanok
    \uses{def:peps_ordered_source_positions,thm:peps_register_identifications}
    Induct on the list of positions. The induction step prepends the actual
    source vector and identifies the unchanged spectator lists.
\end{proof}

\begin{theorem}[Uniform permutation of source registers]
    \label{thm:peps_reorder_source_positions}
    \lean{TNLean.PEPS.PairEffect.SourceInventory.exists_reorderSourcePositions}
    \lean{TNLean.PEPS.PairEffect.SourceInventory.exists_reorderFreeSlots}
    \leanok
    If $J$ and $J'$ are permutations of the same list of source positions, fix
    a spectator register list. There is an allowed composition $U$ containing
    no sources such that, for every source-vector family,
    \begin{align}
        U P_J(\eta) & = P_{J'}(\eta).
        \label{eq:peps_frames_reorder}
    \end{align}
    In particular, one may place any selected subset of free positions before
    its complement, preserving the order within each subset. The same $U$
    works for every vector family; no normalization is required.
\end{theorem}
\begin{proof}
    \leanok
    \uses{thm:peps_ordered_source_preparation}
    Induct on a decomposition of the permutation into adjacent exchanges.
    Use the identity for the empty list, retain a common first pair as
    spectators of the inductively constructed composition, exchange adjacent two-register
    blocks, and compose successive permutations. Apply this result to the
    selected list followed by its complement.
\end{proof}

\begin{definition}[The joint vector of selected sources]
    \label{def:peps_free_source_vector}
    \lean{TNLean.PEPS.PairEffect.SourceInventory.freeSourceVector}
    \lean{TNLean.PEPS.PairEffect.SourceInventory.crossingSourceVector}
    \lean{TNLean.PEPS.PairEffect.SourceInventory.mem_freeSlotLayout}
    \leanok
    For selected positions $J$, let $\eta_J\in\mathcal H_{L_J}$ be their ordered
    tensor product, with the scalar $1$ for the empty list. Its memory is fixed
    by the endpoint spaces, independently of the vectors. Each register is
    exactly one of the two registers at a selected original position.
    For a party division, define the crossing vector by relabelling the
    selected registers by their side and applying $\Pi_{L_J}$ to $\eta_J$.
\end{definition}

\begin{theorem}[Preparation and normalization of source vectors]
    \label{thm:peps_free_source_norm}
    \lean{TNLean.PEPS.PairEffect.SourceInventory.eval_prepareFreeSlots_eq_appendIso_symm}
    \lean{TNLean.PEPS.PairEffect.SourceInventory.norm_freeSourceVector}
    \lean{TNLean.PEPS.PairEffect.SourceInventory.norm_crossingSourceVector}
    \lean{TNLean.PEPS.PairEffect.SourceInventory.freeSourceVector_congr}
    \leanok
    For every spectator vector $x$,
    \begin{align}
        P_J(\eta)x & = \alpha_{L_J,M}^{-1}(\eta_J\otimes x).
        \label{eq:peps_frames_prepare}
    \end{align}
    The joint vector depends only on the selected vectors. If each selected
    vector has norm one, both $\eta_J$ and its crossing vector have norm one,
    including when $J$ is empty.
\end{theorem}
\begin{proof}
    \leanok
    \uses{def:peps_free_source_vector}
    Preparation gives the ordered tensor product by induction. The norm of
    an elementary tensor is the product of the norms; the register
    identifications and the party partition are isometries. Equality on the
    selected positions gives equality of the ordered source records.
\end{proof}

\begin{definition}[Corrected and crossing blocks]
    \label{def:peps_corrected_crossing_blocks}
    \lean{TNLean.PEPS.PairEffect.SourceCircuit.correctedSlotMask}
    \lean{TNLean.PEPS.PairEffect.SourceCircuit.crossingSlotMask}
    \lean{TNLean.PEPS.PairEffect.SourceCircuit.correctedSourceLayout}
    \lean{TNLean.PEPS.PairEffect.SourceCircuit.crossingSourceLayout}
    \leanok
    Let $S$ be a set of original source occurrences of a circuit $W$, and let
    $A(S)$ consist of all their endpoint parties. Among the surviving partial
    source positions, call a position corrected if its original occurrence is
    in $S$, and crossing if its two endpoints lie on opposite sides of $A(S)$.
    Write $C$ and $T$ for the corrected and crossing register lists, respectively,
    with all lists in the inherited original order. Their Boolean-labelled
    versions record only the side of each register.
\end{definition}

\begin{theorem}[The free positions are corrected or crossing]
    \label{thm:peps_corrected_crossing_masks}
    \lean{TNLean.PEPS.PairEffect.SourceCircuit.correctedSlotMask_endpoints}
    \lean{TNLean.PEPS.PairEffect.SourceCircuit.correctedFreeSlots_eq_or}
    \lean{TNLean.PEPS.PairEffect.SourceCircuit.crossingSlotMask_eq_false_of_corrected}
    \lean{TNLean.PEPS.PairEffect.SourceCircuit.correctedFreeSlots_masks}
    \lean{TNLean.PEPS.PairEffect.SourceCircuit.correctedSourceLayout_owners}
    \lean{TNLean.PEPS.PairEffect.SourceCircuit.restrict_correctedSourceLayout_append}
    \lean{TNLean.PEPS.PairEffect.SourceCircuit.correctedMask_meets}
    \leanok
    Both endpoints of every corrected position lie in $A(S)$. Thus corrected
    and crossing positions are disjoint, and the free positions are exactly
    their union. All registers in $C$ lie on the affected side. Consequently,
    the two restrictions of $C+T$ are $C+T_A$ and $T_E$, where $T_A,T_E$ are
    the affected and exterior restrictions of $T$. Every occurrence in $S$
    meets $A(S)$.
\end{theorem}
\begin{proof}
    \leanok
    \uses{def:peps_corrected_crossing_blocks}
    Membership in $A(S)$ is defined by being an endpoint of an occurrence
    in $S$. Apply this to each corrected position. The two Boolean masks are
    then disjoint, and their union is precisely the defining free mask.
\end{proof}

\begin{theorem}[A common corrected-source rearrangement]
    \label{thm:peps_corrected_reordering}
    \lean{TNLean.PEPS.PairEffect.SourceCircuit.exists_correctedInputReordering}
    \leanok
    There is an allowed composition $U$, containing no sources, from the
    corrected block followed by the crossing block and original input to the
    original free-register order. It satisfies~(\ref{eq:peps_frames_reorder})
    for every free-source vector family. In particular, $U$ is chosen before
    the partial branch and before any corrected-source coordinate.
\end{theorem}
\begin{proof}
    \leanok
    \uses{thm:peps_reorder_source_positions,thm:peps_corrected_crossing_masks}
    Apply the uniform permutation theorem to the free positions, selecting
    the corrected positions. The complementary list is the crossing list by
    the disjointness of the two masks.
\end{proof}

\begin{theorem}[Preparing vectors from the scalar unit]
    \label{thm:peps_source_scalar_preparation}
    \lean{TNLean.PEPS.PairEffect.SourceInventory.eval_prepare_nil_heq}
    \lean{TNLean.PEPS.PairEffect.SourceInventory.vector_append_heq}
    \lean{TNLean.PEPS.PairEffect.SourceInventory.eval_prepareFreeSlots_nil_heq}
    \lean{TNLean.PEPS.PairEffect.SourceInventory.vector_append_eq}
    \lean{TNLean.PEPS.PairEffect.SourceInventory.eval_framed_sourceVector}
    \lean{TNLean.PEPS.PairEffect.SourceInventory.eval_framed_freeSourceVectors}
    \leanok
    Preparing an inventory on $1\in\mathbb C$ gives its ordered joint vector.
    Concatenating two inventories gives the vector
    $\alpha_{L,M}^{-1}(\eta_L\otimes\eta_M)$. If a scalar-input composition
    $p_0$ acts with these registers as spectators, it acts as the identity on their
    joint vector and prepares $p_0(1)$ in the original input memory.
\end{theorem}
\begin{proof}
    \leanok
    \uses{thm:peps_free_source_norm,thm:peps_register_identifications}
    Use the preparation identity~(\ref{eq:peps_frames_prepare}) and the
    scalar tensor identity. Tensoring with the identity commutes with preparing the unchanged
    source registers. Apply this once to an inventory and then to its two
    concatenated parts.
\end{proof}

\begin{definition}[Absorbing the original input preparation]
    \label{def:peps_prepared_reordered_word}
    \lean{TNLean.PEPS.PairEffect.SourceInventory.sourceVectorPair}
    \lean{TNLean.PEPS.PairEffect.SourceInventory.preparedReorderedWord}
    \lean{TNLean.PEPS.PairEffect.SourceCircuit.correctedInputPreparation}
    \lean{TNLean.PEPS.PairEffect.SourceCircuit.correctedPreparedWord}
    \lean{TNLean.PEPS.PairEffect.SourceCircuit.correctedInputVector}
    \leanok
    Given a scalar-input composition $p_0$, a register permutation $U$, fixed
    source vectors, and a residual composition $V$, form $\widetilde V$ by
    retaining the two free blocks as spectators of $p_0$, applying $U$, preparing the
    fixed sources, and applying $V$. Its input is just $C+T$.
    The corresponding input vector is
    $\alpha_{C,T}^{-1}(\eta_C\otimes\eta_T)$.
    For a partial circuit branch $\xi$, use its actual residual and its
    original fixed-source vectors; denote the resulting composition by
    $\widetilde V_\xi$.
\end{definition}

\begin{theorem}[Evaluation after absorbing the input]
    \label{thm:peps_prepared_reordered_vector}
    \lean{TNLean.PEPS.PairEffect.SourceInventory.eval_reordered_prepared_vector}
    \lean{TNLean.PEPS.PairEffect.SourceCircuit.correctedInputVector_identity_of_reordering}
    \leanok
    Suppose $U$ has the uniform preparation property
    (\ref{eq:peps_frames_reorder}). Then $\widetilde V$ evaluated on
    $\alpha_{C,T}^{-1}(\eta_C\otimes\eta_T)$ equals the original residual
    applied after full source preparation on $p_0(1)$, with the prescribed
    fixed vectors outside the free positions. The same identity holds for
    every actual partial branch and every free-vector family.
\end{theorem}
\begin{proof}
    \leanok
    \uses{def:peps_prepared_reordered_word,thm:peps_source_scalar_preparation}
    First move the input preparation past the two free-source
    preparations. Apply the defining equality for $U$, then compose free
    preparation with fixed preparation. Their composition is exactly full
    source preparation.
\end{proof}

\begin{definition}[Uniform equality of corrected inputs]
    \label{def:peps_uniform_corrected_input}
    \lean{TNLean.PEPS.PairEffect.SourceCircuit.HasCorrectedInputVectors}
    \leanok
    The uniform corrected-input identity asserts the existence of one allowed,
    source-free $U$ such that the evaluation equality in
    Theorem~\ref{thm:peps_prepared_reordered_vector} holds for every partial
    branch and every free-source family, using the actual residual and fixed
    vectors of that branch.
\end{definition}

\begin{theorem}[Construction of the uniform input identity]
    \label{thm:peps_uniform_corrected_input}
    \lean{TNLean.PEPS.PairEffect.SourceCircuit.exists_correctedInputVector_identity}
    \leanok
    Every circuit and scalar-input composition satisfy the uniform
    corrected-input identity. No normalization or contraction assumption is
    needed for this equality.
\end{theorem}
\begin{proof}
    \leanok
    \uses{thm:peps_corrected_reordering,thm:peps_prepared_reordered_vector}
    Theorem~\ref{thm:peps_corrected_reordering} gives the common permutation.
    Apply Theorem~\ref{thm:peps_prepared_reordered_vector} separately to each
    branch and vector family.
\end{proof}

\begin{theorem}[Actual local factors of the rearranged composition]
    \label{thm:peps_reordered_local_factors}
    \lean{TNLean.PEPS.PairEffect.SourceCircuit.hasTensorPartition_reordered_partial}
    \lean{TNLean.PEPS.PairEffect.SourceCircuit.hasTensorPartition_prepared_reordered_partial}
    \leanok
    Suppose the circuit is allowed. After the common rearrangement, its
    selectively prepared residual factors across the affected/exterior
    division into two actual allowed compositions with no sources and
    operator norms at most one. If $p_0$ is also allowed and contains no
    sources, the same assertion holds after absorbing $p_0$; the free input
    is then exactly $C+T$. The factorization holds on every free input vector.
\end{theorem}
\begin{proof}
    \leanok
    \uses{thm:peps_corrected_crossing_masks,def:peps_prepared_reordered_word}
    Every source remaining in the selectively prepared composition is
    internal to the division. Relabel by the side, so these preparations
    become local operations, and apply the source-free binary factorization.
    Retaining spectators and composing the allowed source-free $p_0$ preserves these
    properties.
\end{proof}

\begin{definition}[Ordered tensor products of source isometries]
    \label{def:peps_source_input_isometries}
    \lean{TNLean.PEPS.PairEffect.SourceInventory.matrixFrame}
    \lean{TNLean.PEPS.PairEffect.SourceInventory.sourcePositionsFrame}
    \lean{TNLean.PEPS.PairEffect.SourceInventory.selectedSourceFrame}
    \leanok
    A matrix $E$ with $E^*E=1$ defines the linear isometry whose columns are
    $E$. Tensor the two endpoint isometries in the order of a list of source
    positions. Restricting this construction to a selected list gives an
    isometry from the Euclidean space indexed by the selected endpoint
    coordinates into its actual register memory. Repeated endpoint parties
    remain distinct tensor factors.
\end{definition}

\begin{theorem}[Columns are the actual prepared endpoint vectors]
    \label{thm:peps_source_input_columns}
    \lean{TNLean.PEPS.PairEffect.SourceInventory.matrixFrame_basis}
    \lean{TNLean.PEPS.PairEffect.SourceInventory.sourcePositionsFrame_basis_prepare}
    \lean{TNLean.PEPS.PairEffect.SourceInventory.selectedSourceFrame_basis}
    \lean{TNLean.PEPS.PairEffect.SourceInventory.selectedSourceFrame_basis_prepare}
    \leanok
    The column at a supplied endpoint assignment is the ordered tensor of its
    literal endpoint columns. With an arbitrary spectator vector appended,
    this column gives exactly preparation of those elementary pair vectors.
    For a selected list, coordinates outside the selected positions make no
    contribution.
\end{theorem}
\begin{proof}
    \leanok
    \uses{def:peps_source_input_isometries,thm:peps_ordered_source_preparation}
    Evaluate each matrix isometry on a coordinate basis vector. Induct on
    the list using the canonical tensor-product basis and source preparation.
    Then reindex the selected coordinate assignments by the filtered list.
\end{proof}

\begin{definition}[The corrected-source isometry]
    \label{def:peps_corrected_source_frame}
    \lean{TNLean.PEPS.PairEffect.SourceCircuit.correctedSourceFrame}
    \leanok
    Fix a partial branch $\xi$ and the original endpoint matrices $E_e,F_e$,
    each with orthonormal columns. For $I=\prod_{e\in S}(J_e\times J_e)$,
    tensor their columns at the actual local labels of $\xi$, in the original
    occurrence order. Relabelling these registers by the affected side gives
    an isometry $Q_\xi\colon\ell^2(I)\to\mathcal H_C$.
\end{definition}

\begin{theorem}[Corrected coordinates preserve the exact source order]
    \label{thm:peps_corrected_source_columns}
    \lean{TNLean.PEPS.PairEffect.SourceCircuit.correctedSourceFrame_basis}
    \leanok
    For every corrected coordinate $i\in I$, $Q_\xi e_i$ is the actual joint
    corrected-source vector formed from the original $E_e,F_e$ columns at $i$.
    The coordinate correspondence is the bijection between selected partial
    positions and original occurrences in $S$.
\end{theorem}
\begin{proof}
    \leanok
    \uses{thm:peps_source_input_columns,thm:peps_free_source_norm}
    Apply the selected-column identity. The existence of the branch and
    the normalization of its original source vectors make each needed
    endpoint coordinate set nonempty. Extend a selected assignment outside
    the selected positions, then remove those choices using dependence only
    on the selected source vectors.
\end{proof}

\begin{definition}[The exact crossing-source vector]
    \label{def:peps_actual_crossing_vector}
    \lean{TNLean.PEPS.PairEffect.SourceCircuit.sourceCrossingVector}
    \leanok
    For branch $\xi$, let $\beta_\xi\in\mathcal H_{T_A}\otimes\mathcal H_{T_E}$
    be the actual ordered tensor of its crossing sources, relabelled and
    partitioned into affected and exterior registers. Its two ambient
    memories depend on the fixed register lists, not on the branch vectors.
\end{definition}

\begin{theorem}[Schmidt isometries for the exact crossing vector]
    \label{thm:peps_crossing_schmidt}
    \lean{TNLean.PEPS.PairEffect.SourceCircuit.norm_sourceCrossingVector}
    \lean{TNLean.PEPS.PairEffect.SourceCircuit.exists_sourceCrossing_schmidt}
    \leanok
    For every partial branch, $\|\beta_\xi\|=1$. There are a finite index set
    $J$, probabilities $\tau_j\geq0$ with $\sum_j\tau_j=1$, and linear
    isometries $L\colon\ell^2(J)\to\mathcal H_{T_A}$ and
    $R\colon\ell^2(J)\to\mathcal H_{T_E}$ such that
    \begin{align}
        \beta_\xi & = \sum_{j\in J}\sqrt{\tau_j}\,Le_j\otimes Re_j.
        \label{eq:peps_frames_crossing_schmidt}
    \end{align}
    The isometries need not span either ambient memory.
\end{theorem}
\begin{proof}
    \leanok
    \uses{thm:peps_free_source_norm,def:peps_actual_crossing_vector}
    Every original source vector is normalized by the prepared-gate construction, so the
    crossing vector has norm one. Place this algebraic tensor in finite
    dimensional subspaces of its two memories, apply the finite dimensional
    Schmidt decomposition, and include the resulting orthonormal families
    back into the original memories.
\end{proof}

\begin{definition}[Schmidt input isometries]
    \label{def:peps_schmidt_input_frames}
    \lean{TNLean.PEPS.PairEffect.HasSchmidtInputFrames}
    \leanok
    Fix the corrected isometry $Q$ and the crossing vector $\beta$.
    Thus $Q\colon\ell^2(I)\to\mathcal H_C$. Choose probabilities $\tau$
    and a linear isometry
    $A\colon\ell^2(I\times J)\to\mathcal H_{(C+T)_A}$.
    Choose also $B\colon\ell^2(J)\to\mathcal H_{(C+T)_E}$.
    A Schmidt input representation requires the identity
    \begin{align}
        \Pi_{C+T}\alpha_{C,T}^{-1}(Qe_i\otimes\Pi_T^{-1}\beta)
         & = \sum_{j\in J}\sqrt{\tau_j}\,Ae_{(i,j)}\otimes Be_j.
        \label{eq:peps_frames_input_schmidt}
    \end{align}
    The probabilities and isometries are chosen before $i$.
\end{definition}

\begin{theorem}[Construction of the actual input isometries]
    \label{thm:peps_actual_schmidt_input_frames}
    \lean{TNLean.PEPS.PairEffect.SourceCircuit.exists_correctedInputFrames}
    \leanok
    For each actual partial branch and original endpoint matrices with
    orthonormal columns, the corrected isometry $Q_\xi$ and exact crossing
    vector $\beta_\xi$ admit the representation
    (\ref{eq:peps_frames_input_schmidt}).
\end{theorem}
\begin{proof}
    \leanok
    \uses{thm:peps_corrected_crossing_masks,thm:peps_crossing_schmidt,def:peps_corrected_source_frame}
    Tensor $Q_\xi$ with the affected isometry in
    (\ref{eq:peps_frames_crossing_schmidt}), and use the exterior isometry
    unchanged. The corrected registers all lie on the affected side, so
    canonical concatenation identifies their targets with the two actual
    restrictions of $C+T$. Distribute the tensor product over the Schmidt
    sum.
\end{proof}

\begin{definition}[A local-word Schmidt output representation]
    \label{def:peps_schmidt_local_output}
    \lean{TNLean.PEPS.PairEffect.Word.HasSchmidtOutput}
    \leanok
    A family $z_i$ in the affected/exterior output tensor product has a
    local-word Schmidt representation if there are actual allowed source-free
    compositions $D_A,D_E$, each of operator norm at most one, probabilities
    $\tau$, and linear input isometries $A,B$ such that
    \begin{align}
        z_i & = \sum_j\sqrt{\tau_j}\,(D_AAe_{(i,j)})\otimes(D_EBe_j).
        \label{eq:peps_frames_output_schmidt}
    \end{align}
    All these choices precede the coordinate $i$.
\end{definition}

\begin{theorem}[Applying the two actual local factors]
    \label{thm:peps_local_schmidt_output}
    \lean{TNLean.PEPS.PairEffect.Word.hasSchmidtOutput_of_hasTensorPartition}
    \leanok
    If a composition factors into actual allowed source-free local words, and
    its inputs admit~(\ref{eq:peps_frames_input_schmidt}), then its outputs
    admit~(\ref{eq:peps_frames_output_schmidt}). The original register
    identifications are retained on both sides.
\end{theorem}
\begin{proof}
    \leanok
    \uses{def:peps_schmidt_input_frames,def:peps_schmidt_local_output}
    Transport the input isometries along the canonical equal-list
    identifications, apply the exact tensor-product factorization to
    (\ref{eq:peps_frames_input_schmidt}), and use linearity. The norm bounds
    come from allowedness of the two local words.
\end{proof}

\begin{definition}[The actual corrected-column input identity]
    \label{def:peps_corrected_schmidt_inputs}
    \lean{TNLean.PEPS.PairEffect.SourceCircuit.HasCorrectedSchmidtInputs}
    \leanok
    The corrected-column input identity asserts that one allowed source-free
    rearrangement, chosen before every branch and corrected coordinate,
    identifies the original partially prepared output with the absorbed
    composition evaluated on the actual corrected-column vector tensor the
    original crossing vector.
\end{definition}

\begin{theorem}[Substitution of the corrected endpoint columns]
    \label{thm:peps_corrected_schmidt_input_identity}
    \lean{TNLean.PEPS.PairEffect.SourceCircuit.exists_correctedSchmidtInput_identity}
    \leanok
    Every circuit, endpoint matrix family and scalar-input composition satisfy
    the corrected-column input identity. No orthonormality or contraction
    hypothesis is needed for this equality.
\end{theorem}
\begin{proof}
    \leanok
    \uses{thm:peps_uniform_corrected_input,thm:peps_corrected_crossing_masks,thm:peps_free_source_norm}
    In the uniform input identity, substitute the original endpoint
    columns at corrected positions and the original source vectors elsewhere.
    Fixed positions are not corrected; crossing positions are not corrected
    either. Dependence only on selected positions therefore preserves the
    literal crossing vector and every fixed source.
\end{proof}

\begin{theorem}[Schmidt expansion of the actual partial output]
    \label{thm:peps_actual_corrected_schmidt_output}
    \lean{TNLean.PEPS.PairEffect.SourceCircuit.exists_correctedSchmidtOutput}
    \leanok
    Let $W$ be an allowed circuit, let all original endpoint matrices have
    orthonormal columns, and let $p_0$ be an allowed scalar-input composition
    containing no sources. For each partial branch $\xi$, replace precisely
    the corrected source vectors by their original endpoint columns at
    $i\in I$, and use $p_0(1)$ as the original input. After the canonical owner
    relabelling and affected/exterior output partition, this actual output
    has the representation~(\ref{eq:peps_frames_output_schmidt}).
    Thus the two local words, probabilities and linear input isometries are
    constructed from the circuit and are chosen before $i$.
\end{theorem}
\begin{proof}
    \leanok
    \uses{thm:peps_corrected_schmidt_input_identity,thm:peps_reordered_local_factors,thm:peps_actual_schmidt_input_frames,thm:peps_local_schmidt_output}
    Choose the common rearrangement from the corrected-column input
    identity. Factor the actual absorbed composition into its two local words.
    Construct its input isometries from the corrected-source isometry and
    exact crossing vector, then apply Theorem~\ref{thm:peps_local_schmidt_output}.
    Finally substitute the original partial-output identity. This keeps the
    original occurrence-coordinate bijection and both owner relabellings.
\end{proof}
